\scanpage{19}
\mspage{\handwritten{10}}

\deleted{\unsure{card(x|card(}}
\[
\operatorname{card}(Y_t\handwritten{(k')}\deleted{\unsure{(t))|t))}}) \equiv a \quad \bmod p\,,
\]
\margin{\handwritten{$|$}}où $Y_t$ est la section de $X_{k\handwritten{'}\deleted{\unsure{(t)}}}$ par l'hypersurface de $P^r_{k\handwritten{'}\deleted{(t)}}$ définie par $t$.

En effet, Deligne a vérifié, à l'aide de la formule de Verdier-Lefschetz, que pour un schéma propre $Y$ sur un corps fini, la somme alternée des traces \deleted{est donc} des frobenius\inserted{géométriques} opérant sur $H^i(Y,\underline{O}_Y)$ est \deleted{\unsure{confru|congru}} égal au nombre des points rationnels de $Y$ mod $p$. D'autre part, quitte à restreindre $U$, nous pouvons supposer que,
\[
f:\underline{Y}\handwritten{\longrightarrow} S_d
\]
étant le morphisme qu'on devine à fibres les $Y_s$ $(s\handwritten{\in}S_d)$, que \inserted{\handwritten{$f$ est plat, et que}} la formation des $R^if_*(\underline{O}_{\underline{Y}})$ commute au changement de base au dessus de $U$, \deleted{et} que les $R^if_*(\underline{O}_Y)$ sont localement libres sur $U$, enfin que $V$ est la fibre \deleted{d'un} générique d'un sous-fibré $\underline{V}$ de $R^{n-1}f_*(\underline{O}_Y)/\underline{O}_{S_d}\otimes H^{n-1}(X,\underline{O}_X)$, stable par frobenius, et sur lequel frobenius soit semi-simple (i.e. un épimorphisme $\underline{V}^{(p)}\handwritten{\longrightarrow}\underline{V}$). Alors la représentation du groupe fondamental du point générique $s$ de $S_d$ envisagée \deleted{e}dans 7.1 provient d'une représentation du groupe fondamental \deleted{d\unsure{e|u}} \handwritten{$\pi$} de $U$, \deleted{\unsure{et même|et ainsi} Quitte à restreindre $U$, on peut supposer également que la partie semi-simple des fibres génériques de tous les $R^if\,(\underline{O}_{\underline{Y}})$ (pas seulement pour $i=n-1$) provient d'un sous-Module localement facteur direct, stable par frobenius, sur lequel frobenius soit semi-simple.} D'autre part, quitte à prendre $d_o$ assez grand dans 7.1, on peut supposer que \deleted{pour} $H^i(X,\underline{O}_X(-d))=0$ pour tout $i\handwritten{<}n$ et tout $d\handwritten{\geq}d_o$ (c'est ici qu'intervient l'hypothèse $X$ Cohen-Macaulay sur $X$), de sorte que pour toute section $Y_s$ de $X$ de degré $d\handwritten{\geq}d_o$ et de dimension $n-1$, l'homomorphisme $H^i(X_s,\underline{\mathbb{O}}_{X_{\handwritten{s}}})\handwritten{\longrightarrow}H^i(\deleted{U}Y_s,\underline{O}_{Y_s})$ sera\handwritten{$|$}un isomorphisme pour $i\leq n-2$, \inserted{\handwritten{un monomorphisme pour $i=n-1$.}} On voit alors que la formule de Lefschetz-Verdier donne, pour \deleted{\unsure{un|la}}\handwritten{le} point \deleted{fermé} $t$ de $U\handwritten{(k')}$, que le nombre des points \deleted{d}rationnel\deleted{\unsure{n|e}}s sur $k\handwritten{'}\deleted{(t)}$ de $Y_t$ \handwritten{satisfait}
% archive sheet number: 18
